% 第1章 研究背景与文献综述
% 证据：tcad-research-training/shared/research-brief.md（导师四项任务与正式题名）；
%       tcad-research-training/shared/literature/mentor-added-20261116.md（L01-L09/J01 拆解）；
%       tcad-research-training/shared/literature/README.md（资料索引与 SHA 清单）；
%       docs/sources/mentor-final-brief.png、mentor-final-title.png（任务与题名截图）；
%       题目_AI辅助下的多特征尺寸存储器单粒子效应评估.md（项目设计大纲提取稿，DOCX SHA256 6ca28c8a…a304）
\section{研究背景与文献综述}

\subsection{研究背景与意义}

动态随机存取存储器（DRAM）以存储电容上的电荷表示信息，访问晶体管导通后存储节点与位线通过电荷共享产生小幅电压差，读出电路据此判断存储状态\cite{keeth2001}。随着特征尺寸持续微缩，存储电容减小、节点电荷减少，单个高能粒子在硅中产生的电子-空穴对经漂移、扩散被节点收集后，就可能改变读出结果，即单粒子效应（SEE）\cite{label1996,baumann2005}。空间环境中的高能质子、重离子以及大气中子是引发 DRAM 单粒子翻转的主要来源。

传统的 SEE 评估高度依赖加速器束流辐照实验：器件筛选、注量累积、错误日志解析都需要机时和样品，成本高、周期长，且难以覆盖全部工艺与粒子工况。因此，将 Geant4 粒子输运仿真与 TCAD 器件仿真联合，并结合数据驱动的预测模型，用有限的实验与仿真数据挖掘多维参数规律、预测翻转截面与错误率，成为降低辐照试验依赖的一条技术路线。这正是导师下发任务的核心动机：调研评估方法与 AI 应用现状，构建 Geant4/TCAD 模型与数据集，建立并训练预测模型，最后与仿真/实验结果对比评价精度与泛化能力\cite{mentortask}。

\subsection{DRAM 单粒子效应评估的研究现状}

\subsubsection{Geant4 与 TCAD 联合仿真}
Geant4 适于描述载能粒子在材料中的输运与能量沉积，TCAD 适于描述器件内部的电荷收集与瞬态电学响应，二者的耦合是 SEE 仿真的主流方法。陈善强等将 TCAD 与 Geant4 用于 SRAM 单粒子效应评估\cite{chen2011}；罗捷等用 Geant4 模拟 Bi、Xe、Kr 的径向沉积并输入 Sentaurus 分析电荷收集，指出 LET 相同也不能保证径向分布相同，粒子种类、能量、位置、时间与通量都应保留在数据链中，其每比特截面定义为 $\sigma_{bit}=N_{bit\ errors}/(\Phi \times N_{tested\ bits})$，分母包含粒子注量与被测比特数\cite{luo2018}。这些工作为本文的 Geant4→TCAD 参数传递与统计分母口径提供了方法依据，但其 SRAM/NMOS 对象与通量阈值不直接照搬到 DRAM。

\subsubsection{器件级建模与校准}
Ryu 等对 14 nm saddle-fin DRAM 开展重离子 TCAD 仿真，给出了结构与 I--V 校准流程，并指出直接收集与邻近节点共享对温度的响应不同，不能用一句"温度越低越抗辐射"概括\cite{ryu2025}；Aditya 等用 14 nm SOI FinFET eDRAM 的 mixed-mode TCAD 分析方向、位置与读出阈值，其报告的临界 LET 依赖特定的读出阈值与电路设置，不是通用值\cite{aditya2018}。这些文献说明：器件模型必须有明确的结构、校准与适用边界，电压下降、共享电荷与最终读错需要分开表述。

\subsubsection{实验表征与错误分类}
Memon 等在 20--58 MeV 质子束下对 COTS LPDDR2/LPDDR4 开展应力测试，按访问图样分组报告计数、注量、截面与不确定度，并强调日志不可解析不等于零错误、SoC 制程与 DRAM 制程不能混用\cite{memon2025}。Jin 等指出 DRAM 的软错误率失效包含阵列位翻转与外围锁存器翻转导致的 SEFI 两类机制，加固比例的分母和层级不同，不能直接当绝对截面\cite{jin2026}。Baumann 的综述区分了 SEU/MBU/SEFI 与节点电荷到系统可见失效的关系\cite{baumann2005}。这些工作为本文的标签统一（瞬态响应、位翻转、SEFI、保持退化分开）与统计分母规范提供了依据。

\subsection{AI 在辐射效应评估中的应用现状}

机器学习已被用于电路级单粒子瞬态的回归建模\cite{xu2021}；Ryder 等比较重离子、聚焦 X 射线与脉冲激光诱导的响应，指出生成电荷总量相近不保证器件瞬态一致，空间分布仍影响响应\cite{ryder2021}。纪志罡的讲义从 DRAM 可靠性的跨层优化视角，讨论了"物理量/波形 → 参数 → 电路指标"的分层代理、图神经网络/迁移学习/PINN 等方法，并强调训练特征应有物理来源、经验拟合好不代表外推可靠\cite{ji2026}。共同趋势是：先建立可解释的简单基线，再比较非线性代理与物理约束模型，并按器件或独立事件分组划分训练/测试，同时报告精度与计算成本。

\subsection{现有不足与本文切入点}

综合上述文献与导师新增资料包，现有工作在以下方面仍有不足：第一，仿真数据集的标签与分母不统一，瞬态响应、判据误读、持久翻转、SEFI 容易混为一谈；第二，Geant4 到 TCAD 的传递与数值验证缺乏可审计的证据链，复跑容易被误当作独立样本；第三，AI 预测多以拟合精度为主，缺少物理约束消融与独立留出的泛化验证；第四，截面/错误率常缺少注量、有效容量与观测时间等分母，难以与实验对照。

本文的切入点是：以可追溯的 Geant4/TCAD 联合仿真链路为基础，统一指标与标签口径，构建带 run ID 与哈希的合格数据集；在此基础上建立数据驱动预测模型并开展超参优化与分组交叉验证，与仿真结果及（待交接的）实验数据对比，分析关键影响因素。

\subsection{研究内容与技术路线}

对应导师任务的四项工作，本文安排为七章：第 2 章介绍 Geant4/TCAD 联合建模与数据链路；第 3 章介绍数据集构建与质量门槛；第 4 章介绍 AI 预测模型、超参优化与交叉验证；第 5 章报告实验结果；第 6 章讨论误差、失败案例与适用边界；第 7 章总结结论与后续工作。技术路线为"已有 DRAM 模型与证据 → 单粒子电学响应和有效读出 → 有统计分母的数据集 → AI 预测与独立验证"。

\subsection{论文组织结构}

第 1 章为背景与文献综述；第 2--4 章为方法（建模、数据集、预测模型）；第 5--6 章为结果与讨论；第 7 章为结论。每章的素材来源、哈希与适用边界统一记录在 \texttt{materials/素材索引.md}，正文定量结论均可回溯。
